\section{Evaluation Details}
\label{app:evaluationdetails}

\subsection{Data and Input Details}
% 前视六个前景类，侧视前四类，均另含背景。
Front-view annotations contain cloth, object, target region, tool, left arm, and right arm; side-view annotations contain the first four classes.
Both views also include background.
% 待核对动作来源、归一化、预处理及真实划分，元信息全90条训练，不可声称未实施的验证。
\draftnote{Specify action targets, normalization, and image preprocessing from the final checkpoints. Verify actual splits and model selection: dataset metadata lists all 90 episodes as training data. Split validation by complete episodes and distinguish acquisition, control, and inference rates.}

\subsection{Visual Completion and Milestone Review}
% 前视主要评价，侧视辅助；完整参考区域不使用被遮挡后的可见掩码。
The front camera is the primary evaluation view; the side camera assists review of occlusion and ambiguous boundaries.
The full target region is defined by an independent reference contour rather than its partially visible segmentation.
% 空掩码不能推断入区，无法判断不确认成功；目标完成时必须仍在区域内。
An empty object mask is not evidence of complete placement, and unobservable intervals are not counted as confirmed success.
The object must remain inside the region at completion.
% 里程碑用于评价，不要求显式阶段输出。
Table~\ref{tab:everydaymilestones} defines evaluation milestones without requiring an explicit phase-prediction head.

\begin{table}[!tb]
\caption{Evaluation milestones for bimanual exposure and execution.}
\label{tab:everydaymilestones}
\centering\small
\begin{tabular}{p{.22\columnwidth}p{.65\columnwidth}}
\toprule
% 里程碑和定义。
Milestone & Definition\\
\midrule
% 暴露：可辨识且达到可见程度。
Exposure & The object becomes identifiable and meets the visibility threshold.\\
% 分离：可见边界不邻接且具有足够可见性。
Separation & Object and cloth boundaries are nonadjacent, with sufficient visibility for assessment.\\
% 入区：完整目标位于参考区域中。
Placement & The entire object lies within the independently defined target region.\\
% 复原：布料画面覆盖达标，目标仍入区。
Restoration & Cloth image coverage reaches the threshold while the object remains fully inside the region.\\
\bottomrule
\end{tabular}
\end{table}
% 累计阶段到达率和条件完成率的分母；零分母不适用，重复不重计。
Let $N_j$ be the number of trials completing the first $j$ milestones in order and $N$ the total number of trials.
Cumulative and conditional completion rates are $N_j/N$ and $N_j/N_{j-1}$, respectively, with $N_0=N$.
Zero denominators are reported as not applicable, and repeated stage entries are not counted again.
% 计时统一，持续确认时记录首达及确认完成两个时刻。
All methods share the same timing origin.
Completion time is the first instant satisfying the full criterion; if persistence is required, use the end of confirmation and also record initial satisfaction.
% 待填写轮廓、边界容差、可见性、邻接、确认时长和审核协议；如容差改变定义须同步正文。
\draftnote{Specify reference contours, boundary tolerance, visibility and adjacency thresholds, confirmation duration, reviewers, disagreement resolution, and unassessable trials. Reflect any allowed placement tolerance in the main criterion.}

\subsection{Trial Allocation and Accounting}
% 六个示范网格位置，基本3类别x12位置，额外未见布料颜色。
Demonstration positions are $(0,0)$, $(0,3)$, $(1,1)$, $(1,2)$, $(2,0)$, and $(2,3)$.
Basic tests cover the three demonstrated objects at all twelve positions; extended tests use an unseen cloth color.
% 相同组合与布置，逐次记录物体位置颜色、checkpoint、种子、重复、时间和失败。
Methods share object--position--color combinations and reset rules.
Each trial records the object, position, cloth color, checkpoint, training seed, repetition, completion time, and primary failure.
% 总体与分组同试验，保留分母，总体不重复相加。
Overall and subgroup results use the same trial pool with explicit denominators.
Overall success counts each trial once, rather than summing overlapping subgroup columns.
% 匹配其他条件，区分位置颜色的独立及联合变化，冻结模型并核查分割开发是否见过。
Position and color comparisons match the remaining conditions and distinguish individual from combined shifts.
Policies and segmenters are frozen before testing, and segmentation-development exposure to the nominally unseen conditions is checked.
% 待定颜色、组合与分母，若分割已见需限定为策略示范未见，尽量保持布料物理特性。
\draftnote{Specify unseen colors, test combinations, and subgroup denominators. Disclose conditions seen during segmentation development; keep cloth material, dimensions, and thickness matched when changing color.}
% 多训练种子，区间保留种子区组；执行重复不等于训练重复。
Key comparisons use independent training seeds, with intervals respecting seed and condition-block structure.
Repeated executions of one checkpoint do not substitute for independent training runs.
% 一次试验一个主要失败类型，其他过程异常单独标签。
Each trial receives one primary failure category; other process anomalies may be recorded separately.
% 待核对清单、样本数、重复和统计方法，36*r仅基础测试，扩展另计；保留全部技术和安全记录。
\draftnote{Verify episode-to-condition mappings, repetitions, seeds, interval estimation, grid spacing, and reset rules. Basic tests contain 36 times the repetitions per condition; count extended tests separately. Prespecify timeouts, safety stops, technical failures, and retry accounting.}

\subsection{Natural Error and Recovery Review}
\label{app:naturalrecovery}
% 所有视频和日志识别自然错误，不只选恢复成功；预先定义可观察判据。
Natural errors are identified from all execution videos and logs, not only successful recovery clips.
The following observable criteria are fixed before evaluation:
\begin{itemize}
% 布料缠阻：持续阻碍；解除并继续任务进展为恢复。
\item \emph{Cloth entanglement:} cloth obstructs tool or object motion; recovery requires clearing the obstruction and resuming task progress.
% 无效接触：推动未接触或无目标进展；重新有效接触并推动为恢复。
\item \emph{Ineffective contact:} an attempted push fails to establish effective contact or object progress; recovery requires subsequent effective contact and displacement.
% 输送偏差：推离路径或停在边缘；重新朝区域输送为恢复，入区另记。
\item \emph{Transport deviation:} the object leaves the transport path or remains at the region boundary; recovery requires renewed transport toward the region, with full placement recorded separately.
% 停滞：任务未完且规定时间无进展；重新有进展为恢复。
\item \emph{Execution stagnation:} no effective progress occurs for a predefined duration before completion; recovery requires renewed progress in the current stage.
\end{itemize}
% 从满足判据到恢复/结束/窗口满；窗口内结束未恢复计失败，恢复需要持续确认。
An event begins when its criterion is first met and ends at recovery, task termination, or expiration of the recovery window.
Recovery requires a confirmation duration; termination without recovery counts as unresolved.
% 连续事件不重复，恢复后再出错才新事件；重叠主要类型，连续偏差停滞不拆分。
Continuous events are counted once; a new event requires recovery followed by recurrence.
Overlapping phenomena receive one primary category and optional auxiliary labels, rather than being split into multiple recovery opportunities.
% 错误发生率按试验，时间仅恢复事件，同时报告分母和出错试验最终成功。
Error incidence counts trials with at least one event.
Recovery time includes recovered events only, alongside recovered/total events and final successes/error-containing trials.
% 无事件不适用，不可审核单列。
Recovery is not applicable when no events occur; unassessable events are reported separately.
% 待定义进展、容差、时长、优先级、窗口，试验聚类统计和严重程度；不完整随访单列，包含失败。
\draftnote{Specify progress criteria, tolerances, stagnation and recovery durations, confirmation, and category precedence. Report severity and counts by condition, use trial-clustered intervals, and list events with incomplete follow-up. Include failed recovery and harmful correction examples.}

\subsection{Representation Diagnostics}
\label{app:evaluationdiagnostics}
% 三组探针：上下文、审核描述量+上下文、预测描述量+上下文；容量/目标匹配，轨迹独立划分。
Expert-action reconstruction probes compare context alone, reviewed descriptors with context, and predicted descriptors with context.
Readout capacity and action targets are matched, with training and evaluation separated by complete trajectories.
% 上下文一致且无未来信息，报告描述量误差、有效比例和动作误差。
Context is identical across probes and excludes target actions and future observations.
Evaluation reports descriptor error, valid-label fractions, and action-reconstruction error.
% 不能替代部署专家标注，也不直接估计覆盖或D。
These probes examine the relation between task information and action selection; they neither replace deployment-side expert labels nor directly estimate coverage or conditional error discrepancy.
% 待填写上下文、动作尺度、读出模型和独立轨迹，速度上一指令需明确因果来源。
\draftnote{Specify context, action scales, readout models, and independent trajectories. Document the causal sources and time ranges of velocity or previous-command inputs.}
